\clearpage\begingroup
\HistoricalHead{191}{R}{WECHSELWIRKUNGEN ZWISCHEN GALVANISCHEN STRÖMEN UND MAGNETISMUS.}
\begin{center}10.\end{center}
\noindent\hspace*{\parindent}Es erhellt daraus, dass der Multiplicator, sobald man den Werth von $v^0$, welcher ihm in Beziehung auf die in ihm befindliche Nadel entspricht, kennt, ein \textls[120]{Galvanometer} im eigentlichen Sinne wird, und die Intensität des ihn durchlaufenden Stroms auf das absolute Maass
\[
g=\frac{T\operatorname{tang}\varphi}{v^0}
\]
\noindent zurückführt. Man erinnere sich dabei, dass $\frac{1}{v^0}$ mit einer Linie homogen ist, z. B. für den oben erwähnten Multiplicator
\[
\frac{1}{v^0}=\frac{1\text{ Millimeter}}{8{,}7196}=0^{\mathrm{mm}}\,1147.
\]
\noindent Dieses ist in Übereinstimmung mit der Bemerkung Art. 3, III [oben S. 180], wonach $g$ mit $RT$ homogen sein muss.

Ein Paar Versuche mit einem hydrogalvanisch erregten Strome mögen eine Idee von der Stärke eines solchen Stromes geben.

I. Ein Paar quadratische Platten von Kupfer und Zink von 5 Zoll Seite durch eine mit Brunnenwasser getränkte Tuchscheibe geschieden, gaben einen Strom, wobei $\varphi=3^0\,27'\,3''$. Hieraus wird
\[
g=0{,}006915\,T,
\]
\noindent oder wenn wir $T=1{,}78$ setzen
\[
g=0{,}012309.
\]
\noindent Hiebei bestand der Leitungsdraht nur aus dem Multiplicatordraht und den kurzen und starken Zuführungsdrähten; dieser Leitungsdraht aequivalirt einem $6070{,}8$ Meter langen Kupferdraht von $1^{\mathrm{mm}}$ Dicke.

II. Indem derselbe Strom zugleich noch die ganze Drahtkette zwischen der Sternwarte, dem magnetischen Observatorium und dem physikalischen Kabinet und die an den beiden letzten Plätzen befindlichen Multiplicatoren durchlief, wobei der ganze Leitungsdraht einem Kupferdraht von $20216$ Meter Länge und $1^{\mathrm{mm}}$ Dicke aequivalirt, war $\varphi=1^0\,32'\,45''$, also
\[
g=0{,}003095\,T=0{,}005509.
\]
\clearpage\SourceLeaf{197}\HistoricalHead{192}{L}{PHYSIK.\quad NACHLASS.}
\noindent Nach dem, was weiter unten entwickelt werden wird, schliesst man hieraus, dass der Widerstand beim Übergang vom Zink durch die Flüssigkeit zum Kupfer so gross war, wie der in einer Kupferdrahtlänge von $5389{,}2$ Meter bei $1^{\mathrm{mm}}$ Dicke, und dass die Stromstärke bei einem kurzen Leitungsdraht von unmerklichem Widerstand
\[
=0{,}014705\,T=0{,}026175
\]
\noindent gewesen sein würde.

Bei dem Werthe $T=1{,}78$ liegen auf die [in] der \textit{Intensitas etc.}[*)]\begingroup\renewcommand{\thefootnote}{}\footnotetext{[*) \textit{Intensitas vis magneticae terrestris ad mensuram absolutam revocata}, 1832, Werke V, S. 79.]}\endgroup\ entwickelte Art als Längeneinheit das Millimeter, als Zeiteinheit die Secunde m. Z. und als Masseneinheit das Milligramm zum Grunde. Will man die drei Einheiten unbestimmt lassen und setzt man sie in ihnen resp. ausgedrückt
\[
\begin{array}{lr}
\text{das Millimeter}&=r\\
\text{die Secunde}&=t\\
\text{das Milligramm}&=p,
\end{array}
\]
\noindent so ist allgemein
\[
T=1{,}78\surd\frac{p}{rtt},
\]
\noindent sowie
\[
v^0=\frac{8{,}7198}{r}.
\]

III. Ein Versuch mit ebenso grossen Platten und demselben Leitungsdraht wie in II, aber unter Anwendung von Tuchscheiben mit gesäuertem Wasser gab
\[
\varphi=1^0\,47'\,40'';\quad g=0{,}003593\,rT=0{,}0063951\surd\frac{pr}{tt}.
\]

IV. Zwei gleiche Plattenpaare als Säule combinirt, unter ähnlichen Umständen wie in III gaben
\[
\varphi=3^0\,22'\,59'';\quad g=0{,}0067795\,rT=0{,}0120675\surd\frac{pr}{tt}.
\]
\noindent Aus Combination dieser beiden Beobachtungen folgt
\begin{flushleft}
der Übergangswiderstand\hfill $=1287{,}8$ Meter,\\[4pt]
die Stromstärke bei kurzem Leitungsdraht\\[2pt]
\hfill $=0{,}059993\,rT=0{,}106788\surd\dfrac{pr}{tt}.$
\end{flushleft}
\clearpage\SourceLeaf{198}\HistoricalHead{193}{R}{WECHSELWIRKUNGEN ZWISCHEN GALVANISCHEN STRÖMEN UND MAGNETISMUS.}
\noindent\hspace*{\parindent}V. Indem der Conductor und das Reibzeug einer Elektrisirmaschine in die Kette gebracht war, wurde die Maschine fortwährend gedreht, 1 Umdrehung auf 1 Secunde. Dabei war die Wirkung $\varphi=1'\,57''\,6$; und der Durchgang der Elektricität aequivalirte einem Strom für welchen
\[
g=0{,}0005701\,rT=0{,}001015\surd\frac{pr}{tt}.
\]
\begin{center}\rule{.18\textwidth}{.35pt}\end{center}
\clearpage\endgroup
